\SourceLeaf{262}
\Article{27}
\indent Si superficies curva est sphaera, cuius radius \(=R\), erit
\[
\begin{gathered}
\alpha=\beta=\gamma=-2f^0=\frac{1}{RR};\quad f''=0,\quad g'=0,\\
6h^0-f^0f^0=0\quad\text{sive}\quad k^0=\frac{1}{24R^4}
\end{gathered}
\]
Hinc formula [14] fit
\[
A+B+C=\pi+\frac{\sigma}{RR}
\]
quae praecisione absoluta gaudet; formulae 11—13 autem suppeditant
\[
\begin{aligned}
A^*&=A-\frac{\sigma}{3RR}-\frac{\sigma}{180R^4}(2pp-qq+4qq'-q'q')\\
B^*&=B-\frac{\sigma}{3RR}+\frac{\sigma}{180R^4}(pp-2qq+2qq'+q'q')\\
C^*&=C-\frac{\sigma}{3RR}+\frac{\sigma}{180R^4}(pp+qq+2qq'-2q'q')
\end{aligned}
\]
sive aeque exacte
\[
\begin{aligned}
A^*&=A-\frac{\sigma}{3RR}-\frac{\sigma}{180R^4}(bb+cc-2aa)\\
B^*&=B-\frac{\sigma}{3RR}-\frac{\sigma}{180R^4}(aa+cc-2bb)\\
C^*&=C-\frac{\sigma}{3RR}-\frac{\sigma}{180R^4}(aa+bb-2cc)
\end{aligned}
\]
Neglectis quantitatibus quarti ordinis, prodit hinc theorema notum a clar. \textsc{Legendre} primo propositum.

\Article{28}
\indent Formulae nostrae generales, reiectis terminis quarti ordinis, persimplices evadunt, scilicet
\[
\begin{aligned}
A^*&=A-\tfrac{1}{12}\sigma(2\alpha+\beta+\gamma)\\
B^*&=B-\tfrac{1}{12}\sigma(\alpha+2\beta+\gamma)\\
C^*&=C-\tfrac{1}{12}\sigma(\alpha+\beta+2\gamma)
\end{aligned}
\]
Angulis itaque \(A,B,C\) in superficie non sphaerica reductiones inaequales applicandae sunt, ut mutatorum sinus lateribus oppositis fiant proportionales. Inaequalitas generaliter loquendo erit tertii ordinis, at si superficies parum a sphaera discrepat, illa ad ordinem altiorem referenda erit: in triangulis vel maximis in superficie telluris, quorum quidem angulos dimetiri licet, differentia semper pro
